% eksperymentalne tłumaczenie na włoski

%

\subsection{Aksjomaty przystawania odcinków} % lang-pl
\index{aksjomat!przystawania odcinków}% % lang-pl
\subsection{Assiomi di congruenza dei segmenti} % lang-it
\index{assioma!di congruenza dei segmenti}% % lang-it

Listę aksjomatów rozszerzymy o trzy, które opisują niezdefiniowane pojęcie przystawania odcinków: relacji między odcinkami, którą oznaczamy przez $\cong$. % lang-pl
(Pompe \cite[s. 25]{pompe_2022} nazywa dwie figury na płaszczyźnie przystającymi, jeśli istnieje izometria, która przekształca jedną na drugą, ale my nie możemy tak zrobić). % lang-pl
Estenderemo la lista degli assiomi con tre nuovi assiomi che descrivono il concetto non definito di congruenza dei segmenti: una relazione tra segmenti che indichiamo con $\cong$. % lang-it
(Pompe \cite[p. 25]{pompe_2022} chiama congruenti due figure del piano se esiste un'isometria che trasforma l'una nell'altra, ma noi non possiamo procedere in questo modo). % lang-it

\begin{axiom}
[przystawania, C1] % lang-pl
[congruenza, C1] % lang-it
    Niech $A$ i $B$ będą punktami, zaś $l$ prostą przechodzącą przez punkt $C$. % lang-pl
    Wtedy po każdej stronie punktu $C$ istnieje punkt $D$ taki, że odcinki $AB \cong CD$ są przystające. % lang-pl
    Siano $A$ e $B$ punti, e sia $l$ una retta passante per il punto $C$. % lang-it
    Allora, da ciascun lato del punto $C$, esiste un punto $D$ tale che i segmenti $AB \cong CD$ siano congruenti. % lang-it
\end{axiom}

Ten aksjomat odpowiada konstrukcji (I.3) Euklidesa i pozwala przenosić odcinki. % lang-pl
\index{odcinek!przenoszenie}% % lang-pl
Questo assioma corrisponde alla costruzione (I.3) di Euclide e permette di trasportare segmenti. % lang-it
\index{segmento!trasporto}% % lang-it

\begin{axiom}
[przystawania, C2] % lang-pl
[congruenza, C2] % lang-it
    Jeśli $AB \cong CD$ oraz $AB \cong EF$, to $CD \cong EF$. % lang-pl
    Każdy odcinek jest przystający do siebie. % lang-pl
    Se $AB \cong CD$ e $AB \cong EF$, allora $CD \cong EF$. % lang-it
    Ogni segmento è congruente a sé stesso. % lang-it
\end{axiom}

To będzie pierwsze pojęcie pierwotne dla Euklidesa. % lang-pl
Questo sarà il primo concetto primitivo per Euclide. % lang-it

\begin{axiom}
[przystawania, C3] % lang-pl
[congruenza, C3] % lang-it
    Niech $A$, $B$, $C$, $D$, $E$, $F$ będą takimi punktami, że $B$ leży między $A$ i $C$, zaś $E$ leży między $D$ i $F$. % lang-pl
    Jeśli $AB \cong DE$ i $BC \cong EF$, to $AC \cong DF$. % lang-pl
    Siano $A$, $B$, $C$, $D$, $E$, $F$ punti tali che $B$ si trovi tra $A$ e $C$, mentre $E$ si trovi tra $D$ e $F$. % lang-it
    Se $AB \cong DE$ e $BC \cong EF$, allora $AC \cong DF$. % lang-it
\end{axiom}

Ten aksjomat pozwala nam dodawać odcinki i znowu zastępuje pojęcie pierwotne Euklidesa. % lang-pl
Patrz też \cite[s. 168]{hartshorne2000}. % lang-pl
Dodawanie jest łączne i przemienne. % lang-pl
Questo assioma ci permette di sommare segmenti e sostituisce nuovamente un concetto primitivo di Euclide. % lang-it
Si veda anche \cite[p. 168]{hartshorne2000}. % lang-it
L'addizione è associativa e commutativa. % lang-it

\begin{proposition}
[odejmowanie odcinków] % lang-pl
[sottrazione di segmenti] % lang-it
    Niech $A$, $B$, $C$, $D$, $E$, $F$ będą takimi punktami, że $B$ leży między $A$ i $C$, zaś $E$ i $F$ leżą na półprostej zaczynającej się w $D$. % lang-pl
    Jeśli $AB \cong DE$ i $AC \cong DF$, to punkt $E$ leży między $D$ i $F$, co więcej $BC \cong EF$. % lang-pl
    Siano $A$, $B$, $C$, $D$, $E$, $F$ punti tali che $B$ si trovi tra $A$ e $C$, mentre $E$ e $F$ appartengano alla semiretta avente origine in $D$. % lang-it
    Se $AB \cong DE$ e $AC \cong DF$, allora il punto $E$ si trova tra $D$ e $F$ e inoltre $BC \cong EF$. % lang-it
\end{proposition}

Odcinek $BC$ traktujemy jako różnicę między $AC$ i $AB$. % lang-pl
Dla Euklidesa to będzie pojęcie pierwotne (że całość jest większa od części). % lang-pl
Consideriamo il segmento $BC$ come la differenza tra $AC$ e $AB$. % lang-it
Per Euclide questo sarà un concetto primitivo (il tutto è maggiore della parte). % lang-it

\begin{definition}
    Niech $AB$ i $CD$ będą odcinkami. % lang-pl
    Jeśli istnieje punkt $E$ między $C$ i $D$ taki, że $AB \cong CE$, to powiemy, że odcinek $AB$ jest krótszy niż odcinek $CD$ (zaś odcinek $CD$ jest dłuższy niż odcinek $AB$). % lang-pl
    Siano $AB$ e $CD$ segmenti. % lang-it
    Se esiste un punto $E$ tra $C$ e $D$ tale che $AB \cong CE$, diremo che il segmento $AB$ è più corto del segmento $CD$ (e che il segmento $CD$ è più lungo del segmento $AB$). % lang-it
\end{definition} % Hartshorne 85

Relacja bycia krótszym odcinkiem jest zgodna z przystawaniem i zadaje zupełny porządek na klasach równoważności. % lang-pl
Dodawanie odcinków zachowuje nierówności między nimi. % lang-pl
La relazione di essere più corto è compatibile con la congruenza e definisce un ordine totale sulle classi di equivalenza. % lang-it
L'addizione dei segmenti preserva le disuguaglianze tra essi. % lang-it

\begin{definition}
[okrąg] % lang-pl
[cerchio] % lang-it
    Niech $O$, $A$ będą dwoma różnymi punktami. % lang-pl
    Zbiór punktów $B$ takich, że $OA \cong OB$ nazywamy okręgiem o środku $O$ oraz promieniu $OA$. % lang-pl
    \index{okrąg} % lang-pl
    Siano $O$ e $A$ due punti distinti. % lang-it
    L'insieme dei punti $B$ tali che $OA \cong OB$ si chiama cerchio di centro $O$ e raggio $OA$. % lang-it
    \index{cerchio} % lang-it
\end{definition}

Okręgi często oznacza się literą $\Gamma$ albo $\Omega$. % lang-pl
Część płaszczyzny ograniczona przez okrąg nazywa się kołem, ale koła nie pojawiają się w części poświęconej planimetrii prawie wcale (poza ustępem o~polu koła). % lang-pl
I cerchi vengono spesso indicati con le lettere $\Gamma$ oppure $\Omega$. % lang-it
La parte di piano delimitata da un cerchio si chiama disco, ma i dischi compaiono molto raramente nella parte dedicata alla planimetria (a eccezione della sezione sull'area del disco). % lang-it

\begin{proposition}
    Każda prosta, która przechodzi przez środek okręgu, przecina go w dwóch punktach. % lang-pl
    Okrąg składa się z nieskończenie wielu punktów. % lang-pl
    Ogni retta che passa per il centro di un cerchio lo interseca in due punti. % lang-it
    Un cerchio è composto da infiniti punti. % lang-it
\end{proposition}

(Nie jest jasne, ile środków może mieć okrąg, ale Hartshorne \cite[s. 89]{hartshorne2000} obiecuje pokazać póżniej, że tylko jeden. % lang-pl
\index{okrąg!środek okręgu}% % lang-pl
Później ma miejsce na stronie 104). % lang-pl
(Non è chiaro quanti centri possa avere un cerchio, ma Hartshorne \cite[p. 89]{hartshorne2000} promette di mostrare più avanti che ne ha uno solo. % lang-it
\index{cerchio!centro del cerchio}% % lang-it
La dimostrazione compare poi a pagina 104). % lang-it

\begin{definition}
[styczna] % lang-pl
[tangente] % lang-it
    Niech $\Gamma$ będzie okręgiem, zaś $l$ prostą, która przecina $\Gamma$ w dokładnie jednym punkcie $A$. % lang-pl
    Mówimy, że $l$ jest styczną do okręgu $\Gamma$ w punkcie $A$. % lang-pl
    \index{styczność} % lang-pl
    Sia $\Gamma$ un cerchio e sia $l$ una retta che interseca $\Gamma$ in un unico punto $A$. % lang-it
    Diremo che $l$ è tangente al cerchio $\Gamma$ nel punto $A$. % lang-it
    \index{tangenza} % lang-it
\end{definition}

Podobnie mówimy, że dwa okręgi są styczne, jeśli mają jeden punkt wspólny. % lang-pl
Neugebauer \cite[s. 20]{neugebauer_2018} podaje w formie ćwiczenia, że prosta i okrąg mają zawsze mniej niż trzy punkty wspólne. % lang-pl
Analogamente, diciamo che due cerchi sono tangenti se hanno un solo punto in comune. % lang-it
Neugebauer \cite[p. 20]{neugebauer_2018} propone come esercizio il fatto che una retta e un cerchio hanno sempre meno di tre punti in comune. % lang-it

%